\chapter{Introducere \c si Motiva\c tie}

\section{Contextul problemei}

Tr\u aim \^ intr-o perioad\u a \^ in care aproape orice aspect al vie\c tii
cotidiene a fost digitalizat -- de la cump\u ar\u aturi la servicii bancare.
Cu toate acestea, \^ in sala de for\c t\u a lucrurile stau surprinz\u ator de
diferit. Mul\c ti sportivi amatori \^ inc\u a \^ i\c si noteaz\u a
antrenamentele \^ intr-un caiet fizic sau, mai r\u au, se bazeaz\u a pe
memorie. Consecin\c ta este mereu aceea\c si: dup\u a dou\u a luni de
antrenament, nu mai \c stii exact ce greut\u a\c ti ai folosit s\u apt\u am\^ ana
trecut\u a \c si nu po\c ti s\u a \^ i\c ti dai seama dac\u a ai progresat sau
ba\c ti pasul pe loc.

Primele aplica\c tii dedicate fitness-ului au ap\u arut \^ in jurul anului
2009, simultan cu explozia smartphone-urilor \cite{grandview2023fitness}.
De atunci, pia\c ta a crescut considerabil, \^ ins\u a direc\c tia de
dezvoltare a majorit\u a\c tii produselor a urmat acela\c si tipar:
func\c tii de baz\u a gratuite, restul \^ in spatele unui abonament lunar.
Rezultatul este c\u a un sportiv care vrea s\u a vad\u a un grafic simplu al
progresului s\u au trebuie s\u a pl\u ateasc\u a, deseori f\u ar\u a s\u a
\c stie exact ce prime\c ste \^ in schimb.

Principiul \textbf{progressive overload} -- ideea c\u a mu\c schii cresc
c\^ and \^ ii supui treptat la efort tot mai mare -- este recunoscut \^ in
literatur\u a ca fundamentul oric\u arui program de antrenament eficient
\cite{zatsiorsky2006science}. Aplicarea lui corect\u a presupune \^ ins\u a
s\u a \c stii \^ intotdeauna de unde ai pornit \c si unde ai ajuns. F\u ar\u a
un istoric clar al antrenamentelor, progressive overload devine aproape
imposibil de aplicat \^ in practic\u a.

\section{Motiva\c tie personal\u a}

Ideea acestui proiect a ap\u arut dintr-o situa\c tie cu care m-am
confruntat personal. Ca student la Informatic\u a care frecventeaz\u a
sala de for\c t\u a de c\u a\c tiva ani, am \^ incercat mai multe aplica\c tii
de tracking -- de la MyFitnessPal la Strong App. Problema era \^ intotdeauna
aceea\c si: fie aplica\c tia \^ imi cerea bani pentru func\c tii pe care le
consider\u a elementare, fie era at\^ at de \^ inc\u arcat\u a de op\c tiuni
\^ inc\^ at m\u a descuraja s-o mai deschid dup\u a antrenament, c\^ and e\c sti
obosit \c si nu ai chef de meniuri complicate.

La un moment dat mi-am zis c\u a, dac\u a tot \^ inv\u a\c t s\u a construiesc
aplica\c tii web, ar fi p\u acat s\u a nu construiesc una care s\u a \^ imi
rezolve exact aceast\u a problem\u a. Nu o aplica\c tie cu zece func\c tii pe
care nu le folose\c ste nimeni, ci una simpl\u a, care s\u a fac\u a bine
un lucru: s\u a \^ imi \c tin\u a minte antrenamentele \c si s\u a \^ imi
arate dac\u a m\u a \^ imbun\u at\u a\c tesc sau nu. Accesibil\u a din orice
browser, gratuit\u a \c si, de ce nu, open-source.

\section{Obiectivele lucr\u arii}

Plec\^ and de la problemele identificate mai sus, aceast\u a lucrare
\^ i\c si propune s\u a realizeze urm\u atoarele:

\begin{enumerate}
  \item Proiectarea \c si implementarea complet\u a a unei aplica\c tii
        web de tip client-server pentru \^ inregistrarea \c si gestionarea
        antrenamentelor de for\c t\u a;
  \item Implementarea unui mecanism de autentificare securizat, cu
        izolarea strict\u a a datelor fiec\u arui utilizator;
  \item Realizarea unui modul de antrenament \^ in timp real, cu suport
        pentru timer de odihn\u a, calculator de distribu\c tie a
        pl\u acilor pe bar\u a \c si indicatori de intensitate
        (RIR \c si 1RM estimat);
  \item Construirea unui modul de analiz\u a a progresului bazat pe
        regresie liniar\u a, cu predic\c tie pe termen scurt \c si
        detec\c tie automat\u a a platourilor de performan\c t\u a;
  \item Validarea solu\c tiei prin utilizare real\u a \c si analiza
        datelor ob\c tinute.
\end{enumerate}

\section{Contribu\c tii originale}

Pe l\^ ang\u a implementarea tehnic\u a propriu-zis\u a, aceast\u a lucrare
aduce urm\u atoarele contribu\c tii originale:

\begin{itemize}
  \item Proiectarea unui model de date adaptat specific domeniului
        antrenamentelor de for\c t\u a, cu suport pentru seturi de
        \^ inc\u alzire, RIR per set \c si asocierea sesiunilor cu
        \c sabloane de rutine;
  \item Implementarea calculului 1RM estimat prin combinarea a dou\u a
        formule consacrate -- Brzycki \cite{brzycki1993strength} \c si
        Epley \cite{epley1985} -- \c si corelarea rezultatului cu
        valoarea RIR \^ inregistrat\u a;
  \item Un algoritm simplu dar eficient de detec\c tie a platourilor,
        bazat pe analiza tendin\c tei volumului s\u apt\u am\^ anal pe
        o fereastr\u a de trei s\u apt\u am\^ ani;
  \item Un calculator de distribu\c tie a pl\u acilor pe bar\u a
        (20 kg), util \^ in practic\u a pentru configurarea rapid\u a
        a greut\u a\c tii la un exerci\c tiu.
\end{itemize}

\section{Structura lucr\u arii}

Lucrarea este \^ imp\u ar\c tit\u a \^ in \c sase capitole, fiecare
construindu-se pe cel anterior.

\textbf{Capitolul 2} analizeaz\u a aplica\c tiile existente pe pia\c t\u a --
MyFitnessPal, Strong, Hevy \c si Jefit -- \c si explic\u a de ce
niciuna nu acoper\u a complet nevoia pe care Gym-Connect o
\^ int\^ ampin\u a.

\textbf{Capitolul 3} prezint\u a toate tehnologiile folosite \^ in
proiect -- de la React \c si TypeScript p\^ an\u a la Prisma,
PostgreSQL pe Neon \c si Clerk -- \c si argumenteaz\u a de ce a
fost aleas\u a fiecare.

\textbf{Capitolul 4} descrie arhitectura sistemului: cum
sunt organizate datele \^ in baza de date, cum func\c tioneaz\u a
API-ul REST \c si cum este gestionat\u a autentificarea.

\textbf{Capitolul 5} prezint\u a aplica\c tia a\c sa cum arat\u a
\^ in realitate -- cu capturi de ecran din fiecare modul \c si
fragmente de cod reprezentative.

\textbf{Capitolul 6} trage concluziile, discut\u a ce ar putea
fi f\u acut mai bine \c si propune direc\c tii de dezvoltare
pentru viitor.
